% Pruebas en teléfono, tercera ronda: la guía con formulario para llenar.
% Se compila con:   cd informes && pdflatex pruebas-telefono-ronda3.tex   (dos veces)
% y después:        python3 -I formulario-radios.py pruebas-telefono-ronda3.pdf pruebas-telefono-ronda3.pdf
% (sin el segundo paso, en algunos visores marcar FALLA no desmarca OK).
% Los campos se leen después con pypdf: resultado-<caso>-<paso> vale ok,
% falla o nopudo; hora-* y nota-* son texto.
\documentclass[10pt,a4paper]{article}

\usepackage[T1]{fontenc}
\usepackage[utf8]{inputenc}
\usepackage[scaled=0.92]{helvet}
\renewcommand{\familydefault}{\sfdefault}
\usepackage[spanish,es-noshorthands]{babel}
\usepackage[a4paper,top=2.2cm,bottom=2.2cm,left=1.9cm,right=1.9cm,headsep=0.8cm]{geometry}
\usepackage{microtype}
\usepackage[table]{xcolor}
\usepackage{array,longtable}
\usepackage{enumitem}
\setlist{itemsep=0.25em, topsep=0.4em, leftmargin=1.2em}
\usepackage{titlesec}
\usepackage{newunicodechar}
\newunicodechar{→}{\ensuremath{\rightarrow}}
\newunicodechar{⋯}{\ensuremath{\cdots}}
\usepackage[most]{tcolorbox}
\usepackage{fancyhdr}
\usepackage{lastpage}
\usepackage[pdfencoding=pdfdoc,hidelinks]{hyperref}
\hypersetup{pdftitle={Pruebas en teléfono: tercera ronda}, pdfauthor={Micros Tempo}}

\setlength{\parindent}{0pt}
\setlength{\parskip}{0.45em}
\renewcommand{\arraystretch}{1.25}
\setlength{\emergencystretch}{2em}

% ── Colores: azul es lo que tocas; marrón, lo que muestra la pantalla ──
\definecolor{azul}{RGB}{28,72,136}
\definecolor{marron}{RGB}{138,78,8}
\definecolor{tenue}{RGB}{105,112,122}
\definecolor{cabecera}{RGB}{233,238,245}
\definecolor{fondoCampo}{RGB}{246,249,253}
\newcommand{\toca}[1]{\textcolor{azul}{\textbf{#1}}}
\newcommand{\pantalla}[1]{\textcolor{marron}{\textbf{#1}}}
\newcommand{\nada}{\textcolor{tenue}{En este paso no hay nada que comprobar.}}

\titleformat{\section}{\large\bfseries\color{azul}}{}{0em}{}
\titlespacing*{\section}{0pt}{1.3em}{0.4em}
\titleformat{\subsection}{\normalsize\bfseries}{}{0em}{}
\titlespacing*{\subsection}{0pt}{0.9em}{0.2em}

\pagestyle{fancy}
\fancyhf{}
\fancyhead[L]{\small\color{tenue}Pruebas en teléfono, tercera ronda · APK 0.3.1 con el arreglo del 20 de noviembre}
\fancyfoot[R]{\small\color{tenue}Página \thepage{} de \pageref{LastPage}}
\renewcommand{\headrulewidth}{0pt}

\newtcolorbox{alto}{colback=red!4, colframe=red!45!black, boxrule=0.5pt, arc=1.5mm,
  left=2.5mm, right=2.5mm, top=1.5mm, bottom=1.5mm, before skip=0.6em, after skip=0.6em}
\newtcolorbox{nota}{colback=azul!4, colframe=azul!55, boxrule=0.5pt, arc=1.5mm,
  left=2.5mm, right=2.5mm, top=1.5mm, bottom=1.5mm, before skip=0.6em, after skip=0.6em}

% ── Los campos del formulario ──────────────────────────────────────
% pdfTeX le da a cada campo el alto de la caja que lo encierra, no el que se
% le pide: sin esta caja, en una fila de la tabla todos medían la fila entera.
\newcommand{\caja}[2]{\hbox{\vrule height #1 depth 0pt width 0pt #2}}
\newcommand{\campo}[3]{% nombre, ancho, alto
  \raisebox{-0.35\height}{\caja{#3}{\TextField[name=#1, width=#2, height=#3, charsize=8pt, bordercolor={0.55 0.6 0.68},
    backgroundcolor={0.965 0.976 0.992}, multiline=true]{}}}}
\newcommand{\dato}[3]{% rótulo, nombre, ancho
  \mbox{#1~\TextField[name=#2, width=#3, height=13pt, charsize=9pt, bordercolor={0.55 0.6 0.68},
    backgroundcolor={0.965 0.976 0.992}]{}}}
\newcommand{\resultado}[1]{%
  \caja{9pt}{\ChoiceMenu[radio, name=resultado-#1, width=9pt, height=9pt, radiosymbol=l,
    bordercolor={0.35 0.38 0.45}, default=]{}{OK=ok,FALLA=falla,NO SE PUDO=nopudo}}}

% Los botones de opción de hyperref ponen la palabra antes del botón, pegada,
% y no declaran su estado inicial (/AS), que algunos visores exigen. Esta copia
% de \@@Radio (hpdftex.def) pone el botón primero, la palabra después con aire,
% y cada botón empieza apagado.
\makeatletter
\def\@@Radio#1{%
  \Fld@listcount=0 %
  \EdefEscapeName\Fld@default{\Fld@default}%
  \@for\@curropt:=#1\do{%
    \expandafter\Fld@checkequals\@curropt==\\%
    \EdefEscapeName\@currValue{\@currValue}%
    \Hy@StepCount\Fld@listcount
    \leavevmode
    \HyAnn@AbsPageLabel
    \Hy@escapeform\PDFForm@Radio
    \pdfstartlink user {%
      \PDFForm@Radio
      /AS/Off
      /AP<</N<</\@currValue\space \DingObject>>>>%
    }%
    \relax
    \MakeRadioField{\Fld@width}{\Fld@height}\pdfendlink
    \ifnum\Fld@listcount=1 \HyField@AddToFields\fi
    \hspace{0.3em}\@currDisplay\hspace{1.1em}%
  }%
}
\makeatother

% Un paso: número visible, número para los campos, qué haces, qué tiene que verse
\newcommand{\paso}[4]{%
  \textbf{#1} & #3 & #4 \\*
  \multicolumn{3}{@{\hspace{6pt}}l@{}}{\rule{0pt}{15pt}\small
    \resultado{#2}\hspace{0.4em}Hora:~\campo{hora-#2}{1.7cm}{13pt}\hspace{1.2em}Nota:~\campo{nota-#2}{8.3cm}{22pt}}\\[3pt]
  \hline
}
\newenvironment{pasos}{%
  \begin{longtable}{@{\hspace{6pt}}>{\raggedright\arraybackslash}p{0.9cm}>{\raggedright\arraybackslash}p{7.9cm}>{\raggedright\arraybackslash}p{7.3cm}@{}}
  \rowcolor{cabecera}\textbf{Paso} & \textbf{Qué haces} & \textbf{Qué tiene que verse} \\ \hline
  \endhead
}{\end{longtable}\vspace{-0.6em}}

\begin{document}
\begin{Form}

{\LARGE\bfseries\color{azul} Pruebas en teléfono: tercera ronda}\par
\vspace{0.2em}
{\color{tenue} APK 0.3.1 con el arreglo del 20 de noviembre de 2026}\par
\vspace{0.6em}

\begin{tabular}{@{}ll@{}}
\dato{Quién prueba:}{tester}{5.6cm} & \dato{Fecha:}{fecha}{3.4cm} \\[4pt]
\dato{Modelo de T1:}{modelo-t1}{5.3cm} & \dato{Modelo de T2:}{modelo-t2}{4.0cm} \\[4pt]
\dato{Chofer A:}{chofer-a}{6.0cm} & \dato{Chofer B:}{chofer-b}{4.6cm} \\[4pt]
\dato{Cobrador de la X:}{cobrador}{4.9cm} & \dato{Combi X:}{combi-x}{2.2cm}\quad\dato{Combi Y:}{combi-y}{2.2cm} \\
\end{tabular}

\section{Para qué es esta ronda}

Esta ronda comprueba en los dos teléfonos los arreglos que se hicieron después de la segunda ronda, y hace por fin el caso del cobrador, que no se pudo hacer en ninguna de las dos anteriores. No hace falta instalar otra APK, porque el arreglo les llega por sí solo a los teléfonos cuando se publica; la preparación dice cómo comprobar que llegó.

Cada caso repite lo que antes falló o no se pudo hacer, y cada paso dice qué haces y qué tiene que verse, en el teléfono y en Despacho. Los casos van en el orden en que conviene hacerlos, que no es el de sus números: el número es el mismo de las rondas anteriores, para que se puedan comparar.

\section{Cómo se anota}

En cada paso marca una de las tres opciones, escribe la hora si algo salió distinto y, si hace falta, una nota corta.

\begin{itemize}
  \item \textbf{OK}: se vio exactamente lo que dice la columna «Qué tiene que verse».
  \item \textbf{FALLA}: se vio otra cosa; anota la hora exacta y saca dos capturas en ese mismo minuto, una del teléfono y otra de Despacho.
  \item \textbf{NO SE PUDO}: el paso no se llegó a hacer; en la nota, escribe por qué.
\end{itemize}

\begin{alto}
\textbf{Si lo que pide la columna «Qué haces» no coincide con lo que ves en la pantalla}, por ejemplo un botón que no está o un texto distinto, detente ahí y anótalo. Desde ese punto, el resto del caso ya no sirve como prueba, aunque después todo parezca andar.
\end{alto}

Para pasarme los resultados, basta con devolverme este PDF lleno. Si me los escribes, usa el número de paso, por ejemplo: «paso 6.4: FALLA, 10:42».

\subsection{Quién es quién}

\begin{itemize}
  \item \textbf{Despacho} es el sistema que se usa en la PC, abierto en el navegador.
  \item \textbf{T1} es el Poco X7 y \textbf{T2} es el Redmi 15C, salvo que anotes otra cosa arriba.
  \item \textbf{El chofer A, el chofer B y el cobrador} son tres usuarios distintos. El cobrador es el que está asignado a la combi X. En la segunda ronda, A era El Gordo, B era Karma, la X era la M11111 y la Y era la M12345.
  \item En las tablas, \toca{lo que está en azul} es lo que tocas o eliges, y \pantalla{lo que está en marrón} es lo que muestra la pantalla. Lo que solo sirve para ubicarte va entre comillas y en negro: eso no se toca.
\end{itemize}

\subsection{Dos cosas que conviene saber antes de empezar}

Para salir de ruta, la app pide dos toques: tocas \toca{SALIR DE RUTA}, el botón cambia a \pantalla{¿SEGURO? TOCA DE NUEVO} y lo tocas otra vez. La guía de la segunda ronda decía que había que mantenerlo apretado, y eso no era exacto.

La combi entra a la ruta cuando el GPS la ve a 20 metros o menos del trazado, más o menos el ancho de la calle. Antes eran 300 metros, unas tres cuadras, y por eso el caso del cobrador no se podía probar sin alejarse cinco cuadras.

% ═══════════════════════════════════════════════════════════════════
\section{Preparación}

Se hace una sola vez, antes de empezar, y deja listos el servidor, los teléfonos y las combis de cada chofer.

\begin{pasos}
\paso{P.1}{P-1}{En la PC, abre Railway y mira el último despliegue del servidor.}{El último despliegue terminó bien y empezó después del 21 de noviembre, cuando el arreglo del 20 entró a \texttt{main}.}
\paso{P.2}{P-2}{En la PC, corre \toca{publicar-arreglo.bat} y elige \toca{[1] PRUEBA}.}{El programa termina sin ningún mensaje que empiece con \pantalla{[ERROR]}.}
\paso{P.3}{P-3}{En cada teléfono, abre la app y espera un minuto en la pantalla de ingreso, sin entrar. Después toca el botón de recientes de Android, desliza la app hacia arriba para cerrarla y ábrela otra vez.}{\nada}
\paso{P.4}{P-4}{En cada teléfono, entra con cualquier usuario, toca \toca{Perfil}, arriba a la derecha, y baja hasta la sección «LA APP».}{Dice \pantalla{Versión 0.3.1 (4) · actualizada} con la fecha y la hora en que publicaste el arreglo. Si no dice \pantalla{actualizada}, el arreglo todavía no llegó: repite el paso P.3.}
\paso{P.5}{P-5}{En Despacho, entra en \toca{Gestión → Personas}, busca la fila del chofer A y toca \toca{⋯} y luego \toca{Combis}. Desmarca la Y y la M-00, que se le habilitaron en la segunda ronda, deja marcada solo la X y toca \toca{Guardar}.}{A queda con una sola combi, la X.}
\paso{P.6}{P-6}{Haz lo mismo con el chofer B, pero deja marcadas dos combis: la X y la Y.}{B queda con dos combis habilitadas.}
\paso{P.7}{P-7}{Deja Despacho abierto en \toca{Mapa} durante toda la prueba. Cuando un paso te mande a la Actividad, vuelve después a \toca{Mapa}. Si algún teléfono tiene una sesión abierta, toca \toca{Salir}, arriba a la derecha.}{Los dos teléfonos muestran la pantalla de ingreso.}
\end{pasos}

% ═══════════════════════════════════════════════════════════════════
\section{Caso 1: el cobrador cuando su combi va hacia la ruta}

T1 lleva al chofer A con la combi X, y T2 lleva al cobrador de la combi X. Este caso se hace en la calle: empiezas a una cuadra o más de la calle por donde pasa el recorrido, y después caminas o manejas hasta pisarla.

\begin{pasos}
\paso{1.1}{1-1}{En T2, entra con el cobrador.}{En T2 aparece \pantalla{ESPERANDO LA RUTA} unos segundos y después \pantalla{NO VEO A TU COMBI}, porque A todavía no salió.}
\paso{1.2}{1-2}{Con T1 a una cuadra o más del recorrido, entra con el chofer A y desliza \toca{DESLIZA PARA SALIR A RUTA →}.}{En T2 aparece \pantalla{TU COMBI VA HACIA LA RUTA}, y se queda fijo: no se va y vuelve.}
\paso{1.3}{1-3}{Espera dos minutos sin moverte.}{En T2 sigue \pantalla{TU COMBI VA HACIA LA RUTA}. Si el aviso se va y vuelve, es FALLA: anota cada cuántos segundos pasa.}
\paso{1.4}{1-4}{Camina o maneja con T1 hasta pisar la calle del recorrido. En la nota, escribe a cuántos metros de la calle estabas cuando cambió la pantalla de T2, más o menos.}{En T2 el aviso se va y aparecen las brechas de la combi X.}
\paso{1.5}{1-5}{Pon T1 en modo avión durante dos minutos.}{En T2 aparece \pantalla{EL CHOFER NO TIENE SEÑAL}, con la hora desde la que no tiene señal.}
\paso{1.6}{1-6}{Quita el modo avión de T1 y no toques nada en T2.}{En T2, el aviso se va solo en menos de un minuto.}
\paso{1.7}{1-7}{En T1, toca \toca{SALIR DE RUTA} y, cuando diga \pantalla{¿SEGURO? TOCA DE NUEVO}, tócalo otra vez.}{En T2 aparece \pantalla{NO VEO A TU COMBI}.}
\paso{1.8}{1-8}{En los dos teléfonos, toca \toca{Salir}, arriba a la derecha.}{Los dos teléfonos muestran la pantalla de ingreso.}
\end{pasos}

% ═══════════════════════════════════════════════════════════════════
\section{Caso 3: AUSENTE cortado por el cierre de sesión}

Este caso se hace en T1, con el chofer A, mirando a la vez el teléfono y Despacho. En la segunda ronda falló una de las dos veces: la combi quedaba en el mapa de Despacho como AUSENTE, sin señal, con la sesión ya cerrada. Por eso ahora se hace tres veces, y cada vez se anota aparte.

\begin{pasos}
\paso{3.1}{3-1}{En T1, entra con el chofer A y desliza \toca{DESLIZA PARA SALIR A RUTA →}.}{Despacho muestra la combi de A en el mapa.}
\paso{3.2}{3-2}{Toca \toca{AUSENTE} y, apenas aparezca \pantalla{Pasas a AUSENTE en 5 s}, toca \toca{Salir}, arriba a la derecha.}{T1 muestra la pantalla de ingreso.}
\paso{3.3}{3-3}{Mira Despacho durante dos minutos, sin tocar nada.}{La combi de A se va del mapa y no vuelve. No debe quedar como \pantalla{AUSENTE} ni como \pantalla{SIN SEÑAL}. Si vuelve, anota a qué hora.}
\paso{3.4}{3-4}{Repite los pasos 3.1, 3.2 y 3.3 por segunda vez.}{Lo mismo que en el paso 3.3: la combi se va del mapa y no vuelve.}
\paso{3.5}{3-5}{Repite los pasos 3.1, 3.2 y 3.3 por tercera vez.}{Lo mismo que en el paso 3.3: la combi se va del mapa y no vuelve.}
\paso{3.6}{3-6}{Entra otra vez con A, sal a ruta y toca \toca{Salir} directamente, sin pasar por AUSENTE. Mira Despacho dos minutos.}{La combi se va del mapa y no vuelve.}
\end{pasos}

% ═══════════════════════════════════════════════════════════════════
\section{Caso 5: cerrar todas las sesiones}

T1 y T2 entran con el mismo usuario, el del chofer A. En este caso no cambies la contraseña: cambiarla también cierra las otras sesiones, y entonces no se sabría si las cerró el botón \toca{CERRAR TODAS MIS SESIONES}.

\begin{pasos}
\paso{5.1}{5-1}{En T1, entra con el chofer A y quédate en la puerta, sin salir a ruta.}{T1 muestra la puerta, con \pantalla{DESLIZA PARA SALIR A RUTA →}.}
\paso{5.2}{5-2}{Cierra la app de T1 desde recientes: toca el botón de recientes de Android y desliza la app hacia arriba.}{La app queda cerrada.}
\paso{5.3}{5-3}{En T2, entra con el chofer A, toca \toca{Perfil} y baja hasta el final. Pasa de largo la sección «CAMBIAR CONTRASEÑA» y su botón «CAMBIAR», sin tocarlos.}{Al final hay un recuadro de borde fino y letra gris que dice \pantalla{CERRAR TODAS MIS SESIONES}.}
\paso{5.4}{5-4}{En T2, toca ese recuadro una vez.}{El recuadro se pone rojo y dice \pantalla{¿SEGURO? TAMBIÉN CIERRA ESTA — TOCA DE NUEVO}.}
\paso{5.5}{5-5}{En T2, tócalo otra vez.}{T2 vuelve a la pantalla de ingreso.}
\paso{5.6}{5-6}{En Despacho, entra en \toca{Gestión → Actividad} y mira la primera línea.}{Dice \pantalla{cerró todas sus sesiones} con el usuario de A, y al lado \pantalla{2 sesión(es)}. Si dice 1, anótalo: quiere decir que una de las dos sesiones ya estaba cerrada.}
\paso{5.7}{5-7}{Espera un minuto, porque el servidor tarda hasta un minuto en cortar las sesiones. Después abre la app en T1.}{Aparece la pantalla de ingreso con el aviso \pantalla{Tu sesión ya no vale: entra de nuevo}. Antes salía «Sesión inválida o expirada», que ya no tiene que aparecer.}
\end{pasos}

\subsection{Variante con la ruta empezada}

Se hace después de los pasos 5.1 a 5.7, y comprueba que el teléfono que estaba en ruta deja de mandar su posición.

\begin{pasos}
\paso{5.8}{5-8}{En T1, entra con el chofer A, desliza \toca{DESLIZA PARA SALIR A RUTA →} y apaga la pantalla con el botón de encendido.}{Despacho muestra la combi de A en el mapa.}
\paso{5.9}{5-9}{En T2, entra con el chofer A, abre \toca{Perfil}, baja hasta el final y toca dos veces \toca{CERRAR TODAS MIS SESIONES}.}{T2 vuelve a la pantalla de ingreso.}
\paso{5.10}{5-10}{Mira Despacho durante dos minutos y anota a qué hora deja de verse la combi de A.}{La combi de A se va del mapa en pocos segundos y no vuelve, aunque T1 siga en ruta con la pantalla apagada: cerrar todas las sesiones también baja al chofer A de su combi en el servidor.}
\paso{5.11}{5-11}{En T1, prende la pantalla y abre la app.}{Aparece la pantalla de ingreso, con el aviso \pantalla{Tu sesión ya no vale: entra de nuevo}. En la barra de notificaciones ya no está la del GPS.}
\end{pasos}

% ═══════════════════════════════════════════════════════════════════
\section{Caso 6: el relevo de chofer en la misma combi}

Este caso necesita los pasos P.5 y P.6 de la preparación: el chofer A tiene una sola combi, la X, y el chofer B tiene dos, la X y la Y. T1 lleva al chofer A y T2 al chofer B. En la segunda ronda, A tomó la combi que manejaba B sin que la app le preguntara nada, porque con una sola combi no aparece la lista; ahora la app tiene que preguntar al deslizar.

\begin{pasos}
\paso{6.1}{6-1}{En T2, entra con el chofer B. En la lista \pantalla{¿EN QUÉ COMBI SALES?}, toca la combi X.}{Las dos combis dicen \pantalla{libre}. La X queda marcada como \pantalla{ELEGIDA}, sin pedir confirmación, porque nadie la maneja.}
\paso{6.2}{6-2}{En T2, desliza \toca{DESLIZA PARA SALIR A RUTA →}.}{T2 queda en ruta. Despacho muestra la combi X con el nombre de B.}
\paso{6.3}{6-3}{En T1, entra con el chofer A y mira la puerta, encima del botón para salir a ruta.}{No hay lista, porque A tiene una sola combi. Encima del botón dice \pantalla{X: la maneja B}, con el código de la X y el nombre de B más grueso que el resto.}
\paso{6.4}{6-4}{En T1, desliza \toca{DESLIZA PARA SALIR A RUTA →}.}{T1 todavía no sale a ruta. Aparece un recuadro que dice \pantalla{X la está manejando B. Si es un relevo, la tomas tú y a B se le avisa.}, con los botones \pantalla{ES UN RELEVO} y \pantalla{Cancelar}.}
\paso{6.5}{6-5}{En T1, toca \toca{Cancelar}.}{El recuadro se cierra y T1 sigue en la puerta. Despacho sigue mostrando la X con B.}
\paso{6.6}{6-6}{En T1, desliza otra vez y, en el recuadro, toca \toca{ES UN RELEVO}.}{T1 queda en ruta con la combi X.}
\paso{6.7}{6-7}{Mira la pantalla de T2 sin tocar nada y saca una captura.}{Aparece el aviso \pantalla{El GPS se apagó: Otro chofer tomó la unidad. Sigues viendo todo, pero tu GPS ya no se usa. Para volver a manejarla: SALIR DE RUTA y después SALIR A RUTA.}}
\paso{6.8}{6-8}{En Despacho, mira la combi X en el mapa y saca una captura. Después entra en \toca{Gestión → Actividad}.}{La X aparece con A, en un solo punto y sin saltos. En la Actividad hay una línea que dice \pantalla{tomó la combi que manejaba}.}
\paso{6.9}{6-9}{En T2, toca \toca{SALIR DE RUTA} dos veces, como en el paso 1.7, y mira la lista de combis de B.}{En la fila de la X dice \pantalla{la maneja A}, con el nombre de A más grueso que el resto.}
\paso{6.10}{6-10}{En T2, toca la fila de la X y, en el recuadro, toca \toca{ES UN RELEVO}.}{La X queda marcada como \pantalla{ELEGIDA} para B.}
\paso{6.11}{6-11}{En T2, desliza \toca{DESLIZA PARA SALIR A RUTA →}.}{T2 sale a ruta sin preguntar otra vez, porque el relevo ya se confirmó en la lista. En T1 aparece el aviso del paso 6.7, ahora para A, y Despacho muestra la X con B.}
\end{pasos}

\subsection{Variante con la pantalla apagada}

Se hace después de los pasos 6.1 a 6.11 y repite el paso que falló en la segunda ronda: B maneja la X con la pantalla apagada y A le hace el relevo.

\begin{pasos}
\paso{6.12}{6-12}{En T1, toca \toca{SALIR DE RUTA} dos veces.}{T1 vuelve a la puerta.}
\paso{6.13}{6-13}{Apaga la pantalla de T2, que sigue en ruta con la X, y anota la hora. Espera dos minutos.}{\nada}
\paso{6.14}{6-14}{En T1, mira la puerta, encima del botón para salir a ruta.}{Sigue diciendo \pantalla{X: la maneja B}. Si dice otra cosa, anota la hora y mira en Despacho si la X aparece como \pantalla{SIN SEÑAL}.}
\paso{6.15}{6-15}{En T1, desliza \toca{DESLIZA PARA SALIR A RUTA →} y, en el recuadro, toca \toca{ES UN RELEVO}.}{T1 queda en ruta con la X. Despacho muestra la X con A, en un solo punto.}
\paso{6.16}{6-16}{En T2, prende la pantalla, abre la app y saca una captura.}{En T2 aparece otra vez el aviso del paso 6.7, para B. En Despacho, la X sigue con A y no vuelve a saltar hacia B.}
\paso{6.17}{6-17}{En los dos teléfonos, sal de ruta con dos toques y toca \toca{Salir}.}{Los dos teléfonos muestran la pantalla de ingreso.}
\end{pasos}

% ═══════════════════════════════════════════════════════════════════
\section{Caso 2: las notificaciones quitadas en ruta}

Este caso se hace en T1, con el chofer A, y va al final porque Android puede cerrar la app a la fuerza. En la segunda ronda, al quitar las notificaciones en ruta, el chofer desapareció de Despacho y no volvió aunque se prendieron otra vez las notificaciones. El arreglo ataca la causa más probable, que no se pudo confirmar sin el teléfono; si vuelve a fallar, lo que más sirve es la captura del paso 2.5.

\begin{pasos}
\paso{2.1}{2-1}{En T1, entra con el chofer A y sal a ruta.}{Despacho muestra la combi de A en el mapa.}
\paso{2.2}{2-2}{En los ajustes de Android, entra en las aplicaciones, busca \toca{Control de ruta}, abre \toca{Notificaciones} y apágalas.}{\nada}
\paso{2.3}{2-3}{Vuelve a la app. Si Android la reinició, entra otra vez con A. En la nota, escribe si la app se reinició sola y si volvió a pedir el permiso de notificaciones.}{\nada}
\paso{2.4}{2-4}{Deja la app abierta y en primer plano durante dos minutos, y mira Despacho.}{Despacho vuelve a mostrar la combi de A en el mapa y sigue recibiendo sus posiciones, porque la app vuelve a encender el GPS sola si en un minuto no le llega ninguna.}
\paso{2.5}{2-5}{Solo si en el paso 2.4 la combi no volvió: toca \toca{Perfil}, baja hasta la sección «EL GPS» y saca una captura.}{La captura muestra qué le pasa al GPS. Si todo salió bien en el paso 2.4, marca NO SE PUDO y escribe «no hizo falta».}
\paso{2.6}{2-6}{En los ajustes de Android, prende otra vez las notificaciones de \toca{Control de ruta} y vuelve a la app sin cerrarla.}{En menos de un minuto vuelve la notificación de la brecha, y Despacho sigue mostrando la combi de A.}
\paso{2.7}{2-7}{Toca \toca{SALIR DE RUTA} dos veces y después \toca{Salir}.}{T1 muestra la pantalla de ingreso, y la combi de A se va del mapa de Despacho.}
\end{pasos}

% ═══════════════════════════════════════════════════════════════════
\section{Al terminar}

Devuélveme este PDF lleno, junto con las capturas de los pasos 6.7, 6.8 y 6.16, las de cada FALLA y, si hizo falta, la del paso 2.5, y yo paso los resultados a \texttt{RESULTADOS-TELEFONO.md}.

\begin{nota}
\textbf{Último paso:} en Despacho, quítale al chofer B la combi X, que le habilitaste en el paso P.6 solo para esta prueba.
\end{nota}

\subsection{Observaciones generales}

Anota aquí lo que no cupo en las notas de cada paso.

\TextField[name=observaciones, width=\linewidth, height=6cm, charsize=9pt, multiline=true,
  bordercolor={0.55 0.6 0.68}, backgroundcolor={0.965 0.976 0.992}]{}

\end{Form}
\end{document}
